\section{Minacce alla validità}

\subsection{Milestone 1}

\paragraph{Linkage commit--ticket}
Il linkage avviene tramite la presenza della chiave del ticket nel messaggio di commit.
Questa assunzione introduce due rischi simmetrici:
commit che correggono un bug senza citare la chiave non vengono associati (falsi negativi),
mentre commit che citano la chiave per ragioni diverse dalla correzione vengono inclusi
erroneamente (falsi positivi).

\paragraph{Stima con Proportion}
Proportion introduce inevitabilmente rumore nel dataset. Tuttavia, l'approccio adottato
dà priorità alle AV esplicite, limitando il ricorso a Proportion ai soli ticket privi di
esse.

\paragraph{Release senza tag Git}
Tre versioni Jira non dispongono di un tag Git corrispondente. I loro snapshot non sono
inclusi nel dataset, il che può comportare una sottostima della bugginess per le classi
che avrebbero dovuto comparire in quelle release.

\paragraph{Utilizzo delle Fix Version come Affected Version}
Per i ticket privi di AV ma con FV disponibili, le FV tranne l'ultima vengono
interpretate come AV. Questa scelta introduce una minaccia alla validità. È quindi possibile
che alcune versioni vengano etichettate come affette da un bug quando in realtà non lo
erano. La scelta è stata comunque preferita al ricorso sistematico a Proportion, che per definizione
costituisce anch'essa una stima e introduce un diverso tipo di rumore.

\subsection{Milestone 2}

\paragraph{Temporal leakage strutturale}
La 10$\times$10 fold CV assegna le istanze ai fold in modo casuale, senza rispettare
l'ordine temporale delle release. Il modello può essere addestrato su informazioni
``future'' rispetto a parte del set di test: nelle versioni più recenti i bug esistenti
nelle versioni precedenti sono già stati corretti, e il modello apprende pattern che non
sarebbero stati disponibili al momento della predizione reale. Le metriche risultanti
tendono a essere leggermente ottimistiche rispetto all'accuratezza in produzione.
Alternativa: Walk-Forward Validation (discussa nella sezione sulla metodologia).

\paragraph{Nessun tuning degli iperparametri}
I classificatori usano i valori di default di Weka (100 alberi per RandomForest, $k=1$
per IBk). Un tuning via grid search migliorerebbe i risultati ma richiederebbe un
ulteriore livello di cross-validation annidata, aumentando ulteriormente i tempi di
esecuzione.

\paragraph{Classificatore interno dei metodi wrapper}
I metodi wrapper usano NaiveBayes come classificatore interno per valutare i subset
candidati. Il subset di feature ottimale per NaiveBayes potrebbe non essere ottimale per
RandomForest o IBk, che vengono poi valutati esternamente con quel subset.

\paragraph{Approssimazione della percentuale SMOTE}
La percentuale di oversampling SMOTE è calcolata sull'intero dataset. La percentuale
esatta per-fold varia del $\approx$10\% (ogni training fold ha il 90\% delle istanze),
introducendo una piccola imprecisione nel bilanciamento.

\paragraph{Limitazione dello spazio di ricerca nei metodi wrapper}
Con 19 feature, GreedyStepwise è computazionalmente trattabile. Con un numero
significativamente maggiore di feature il costo crescerebbe quadraticamente e la ricerca
greedy diventerebbe meno affidabile nel trovare l'ottimo globale.

\subsection{Milestone 3}

\paragraph{Training bias su A}
BClassifierA è addestrato sull'intero dataset A e poi usato per predire A stesso.
Le predizioni su A non sono quindi ``out-of-sample'' e tendono a essere ottimistiche.
Le predizioni su B+, B e C sono invece genuine (A è il training set, B+/B/C sono
applicazioni del modello a dati che non ha mai visto come test).

\paragraph{Indipendenza tra NSmells e le altre feature}
L'analisi controfattuale assume implicitamente che sia possibile azzerare NSmells
mantenendo invariate le altre 18 feature. Nella realtà questo non è realizzabile:
un file privo di smell ha tipicamente anche meno LOC, meno storia di modifica e
maggiore semplicità strutturale. Azzerare solo NSmells nel dataset costruisce un
controfattuale che non corrisponde a nessuna distribuzione reale di classi Java.

\paragraph{Composizione dei ruleset PMD e correlazione con LOC}
La metrica \texttt{Previous\_Release\_Code\_Smells} aggrega le violazioni di quattro
ruleset PMD: \texttt{design}, \texttt{errorprone}, \texttt{bestpractices} e
\texttt{codestyle}. Il ruleset \texttt{codestyle} rileva violazioni di stile puro
(naming convention, stile delle parentesi, import superflui) le cui occorrenze crescono
linearmente con la dimensione del file. Questo introduce una correlazione artificiale
tra NSmells e LOC, che a sua volta riduce l'indipendenza di NSmells rispetto alle
metriche di processo cumulative (LOC\_touched, Churn, NR). La conseguenza è che
RandomForest non ha bisogno di NSmells come feature separata, perché il suo segnale
è già assorbito dalle metriche di processo. Rimuovere \texttt{codestyle} e aggiungere
ruleset meno dimensione-dipendenti (\texttt{performance}, \texttt{security}) potrebbe
ridurre questa correlazione e aumentare il contributo indipendente di NSmells.
Studi diversi usano metriche di smell diverse (spesso SonarQube con smell architetturali
quali God Class, Feature Envy, Long Method), il che limita ulteriormente la
comparabilità dei risultati con la letteratura.

\paragraph{Dipendenza dal classificatore scelto}
I risultati dell'analisi controfattuale dipendono da quale classificatore viene scelto
come BClassifier. Un classificatore diverso (es.\ IBk) avrebbe pesi diversi per
NSmells e potrebbe produrre un numero di istanze prevenute significativamente diverso.

\subsection{Milestone 4}

% TODO: descrivere i limiti del refactoring automatico, le assunzioni sui test
% e le possibili conseguenze di scelte soggettive nel prompt.
